% =============================================================================
% ContextOS -- Cover Letter to Information Sciences
% =============================================================================
% Compile with: pdflatex cover_letter.tex
% =============================================================================

\documentclass[12pt, a4paper]{letter}

\usepackage[T1]{fontenc}
\usepackage[utf8]{inputenc}
\usepackage{microtype}
\usepackage{geometry}
\geometry{
  top=2.5cm,
  bottom=2.5cm,
  left=3.0cm,
  right=3.0cm
}
\usepackage{hyperref}
\hypersetup{
  colorlinks=true,
  linkcolor=black,
  urlcolor=black,
  citecolor=black
}
\usepackage{parskip}
\setlength{\parskip}{0.8em}
\setlength{\parindent}{0pt}

% Signature block spacing
\longindentation=0pt

\begin{document}

\begin{letter}{
  Editor-in-Chief \\
  \textit{Information Sciences} \\
  ISSN: 0020-0255 \\
  Elsevier B.V. \\
  Radarweg 29, 1043 NX Amsterdam \\
  The Netherlands
}

\opening{Dear Editor-in-Chief,}

We are pleased to submit our manuscript titled:

\begin{center}
\textbf{ContextOS: An Operating System for Context Lifecycle Management\\
in Long-Horizon Autonomous Agents}
\end{center}

for consideration as a \textbf{Regular Article} in \textit{Information
Sciences}.  The manuscript is approximately 35 pages in length (including
references and appendices) and has not been previously published, nor is it
under review at any other journal or conference.  All authors have approved
the submission and declare no conflicts of interest.

\textbf{Summary of contributions.}  Autonomous AI agents operating over
long-horizon tasks are fundamentally constrained by the finite capacity of
their context windows.  Existing approaches -- including naive truncation,
retrieval-augmented generation (RAG), and paged memory systems such as MemGPT
-- treat context management as a secondary concern rather than a core
architectural primitive.  This paper introduces \textbf{ContextOS}, an
open-source framework that reframes context as a first-class system resource,
analogous to how operating systems manage CPU time, RAM, and persistent
storage.

ContextOS provides a unified pipeline comprising six coordinated subsystems:
(1) a \emph{Retrieval Engine} for semantic access to episodic, semantic, and
procedural memory; (2) a \emph{Prioritization Engine} that scores context
items using a multi-signal formula combining relevance, recency, importance,
and novelty; (3) a priority-based \emph{Context Scheduler} with greedy,
MMR-diversity, and threshold strategies; (4) a multi-strategy
\emph{Compression Engine} supporting extractive, abstractive, hierarchical,
and progressive compression; (5) a \emph{Working Memory} module for
short-term agent state; and (6) a \emph{Governance Engine} implementing
configurable retention, forgetting, and promotion policies with exponential
decay and LRU eviction.

We introduce \textbf{ContextBench}, a programmatically generated benchmark
dataset of 300{,}000 labelled examples spanning five subsets (multi-task
context management, retrieval, compression, embedding fine-tuning, and LLM
instruction tuning) across six professional domains and eight task types.
ContextBench is fully reproducible from fixed random seeds with no external
dependencies.

Our empirical evaluation on ContextBench, LongBench, MuSiQue, NarrativeQA,
and HotpotQA demonstrates that ContextOS achieves statistically significant
improvements over all baselines across all context lengths.  At 8{,}192
tokens, ContextOS improves Task Success Rate (TSR) by $+7.6$ percentage
points over MemGPT (the strongest prior system) with $p < 10^{-5}$
(Bonferroni-corrected).  At 32{,}768 tokens, the improvement grows to
$+12.4$ points.  Token efficiency is improved by 40.4\% compared to
full-context baselines.

\textbf{Relevance to \textit{Information Sciences}.}  This work falls squarely
within the scope of \textit{Information Sciences}, which publishes research on
intelligent information processing, knowledge engineering, and AI systems.
ContextOS contributes to all three dimensions: (i) principled information
selection and scheduling for AI agents; (ii) knowledge representation across
episodic, semantic, and procedural memory tiers; and (iii) the design and
empirical evaluation of an end-to-end AI system for long-horizon task
completion.  The paper includes rigorous experimental methodology, a
large-scale benchmark dataset, and complete open-source code, making it a
substantial and reproducible contribution to the community.

\textbf{Suggested reviewers.}  We respectfully suggest the following three
experts in context management, agent AI, and information retrieval as
potential reviewers.  We have no conflicts of interest with any of them.

\begin{enumerate}
  \item \textbf{Prof.\ Danqi Chen} (Princeton University) -- Expert in dense
        retrieval, open-domain question answering, and large language model
        contexts.  Her work on REALM, DPR, and FiD directly intersects with
        the retrieval and long-context aspects of this paper.

  \item \textbf{Dr.\ Shunyu Yao} (Princeton University / OpenAI) -- Known for
        ReAct, Tree of Thoughts, and agent reasoning frameworks.  His expertise
        in agent planning and memory management makes him well-placed to
        evaluate ContextOS's scheduling and governance contributions.

  \item \textbf{Prof.\ Heng Ji} (University of Illinois Urbana-Champaign) --
        Expert in information extraction, knowledge graphs, and multi-document
        summarisation.  Her research on information fusion and cross-document
        reasoning is closely related to ContextOS's multi-signal prioritization
        and compression components.
\end{enumerate}

\textbf{Data and code availability.}  The complete ContextOS codebase,
ContextBench dataset generation scripts, and all experimental result files are
available as open source under the Apache License 2.0.  The repository URL
will be made publicly accessible upon acceptance; during review, the code is
available at the anonymised repository link provided in the supplementary
materials.

\textbf{Ethics statement.}  All datasets used in this work are either
synthetic (ContextBench, generated programmatically with no human subjects) or
publicly available benchmarks (LongBench, MuSiQue, NarrativeQA, HotpotQA).
No personal or sensitive data is collected, processed, or published.

We sincerely hope that this submission will be of interest to the readership of
\textit{Information Sciences}.  We are committed to addressing any reviewer
concerns promptly and thoroughly.

\closing{Respectfully yours,}

\vspace{1.5cm}

\noindent
\textit{(Author identities withheld for double-blind review)}

\vspace{0.8cm}

\noindent
\textbf{Corresponding Author:} Anonymous (identities revealed upon acceptance)\\
\textbf{Email:} \texttt{[redacted for review]}\\
\textbf{Affiliation:} \texttt{[redacted for review]}\\
\textbf{Manuscript type:} Regular Article\\
\textbf{Estimated length:} $\approx$35 pages (main text + appendices)\\
\textbf{Number of figures:} 9\\
\textbf{Number of tables:} 12 (main) + 17 (appendix)\\
\textbf{Supplementary material:} Appendix (pseudocode, extended results,
hyperparameter analysis, implementation details)\\
\textbf{Open-source code:} Anonymised repository link in supplementary\\
\textbf{Dataset availability:} Fully reproducible from provided scripts\\
\textbf{Conflicts of interest:} None

\end{letter}

\end{document}
